\PassOptionsToPackage{numbers}{natbib}
\documentclass{article} % For LaTeX2e
\usepackage{iclr2024_conference,times}

\usepackage[utf8]{inputenc} % allow utf-8 input
\usepackage[T1]{fontenc}    % use 8-bit T1 fonts
\usepackage{hyperref}       % hyperlinks
\usepackage{url}            % simple URL typesetting
\usepackage{booktabs}       % professional-quality tables
\usepackage{amsfonts}       % blackboard math symbols
\usepackage{nicefrac}       % compact symbols for 1/2, etc.
\usepackage{microtype}      % microtypography
\usepackage{titletoc}

\usepackage{subcaption}
\usepackage{graphicx}
\usepackage{amsmath}
\usepackage{multirow}
\usepackage{color}
\usepackage{colortbl}
\usepackage{cleveref}
\usepackage{algorithm}
\usepackage{algorithmicx}
\usepackage{algpseudocode}
\usepackage{tikz}
\usepackage{pgfplots}
\usepackage{float}
\usepackage{array}
\usepackage{tabularx}
\pgfplotsset{compat=newest}


\DeclareMathOperator*{\argmin}{arg\,min}
\DeclareMathOperator*{\argmax}{arg\,max}

\graphicspath{{../}} % To reference your generated figures, see below.

\title{Regulator-Ready Worst-Group Strategy Selection with Stratified Multi-Proxy Prediction-Powered E-Audits}

\author{AIRAS}

\newcommand{\fix}{\marginpar{FIX}}
\newcommand{\new}{\marginpar{NEW}}

\begin{document}

\maketitle

\begin{abstract}
Deployment decisions for large language model (LLM) “reasoning strategies” increasingly resemble regulatory approvals: a strategy should be authorized only if it clears a safety or legality threshold for every protected subgroup, and among those group-safe strategies we should deploy the most useful one. This is hard because gold audits are expensive, cheap automated judges are biased in subgroup-specific ways, and practical auditing is adaptive with optional stopping, making common confidence estimates unreliable. We propose Stratified, Multi-Proxy Prediction-Powered E-Audits (SMP-PPeA), an anytime-valid and distribution-free protocol that combines subgroup-stratified proxy scores available for all items with a small stream of gold audits. SMP-PPeA maintains stratified lower confidence bounds using e-process inversion, audits only proxy–gold correction terms, and hedges across multiple proxies per subgroup via conservative risk allocation. We implement a fully scriptable Python pipeline on RealToxicityPrompts with deterministic legality checks and ROUGE-L utility, simulate adaptive auditing across several gold budgets, and evaluate two open-weight models. The executed runs expose a practically important degenerate regime: at a strict worst-group threshold, no strategy is feasible under full gold evaluation, so all methods abstain; we analyze what this implies for designing regulator-ready benchmarks and method comparisons.
\end{abstract}

\section{Introduction}
\label{sec:intro}
Regulators and deployers increasingly expect that decisions about deploying large language model (LLM) systems be supported by auditable evidence rather than point estimates. This is especially salient for legality and safety constraints, where being “safe on average” is not an acceptable substitute for being safe for every protected subgroup. In many real deployments, a single base model can be operated under multiple “reasoning strategies,” such as different instruction prefixes, prompting styles, or decoding policies. Governance then reduces to a high-stakes selection problem: choose which strategy is permitted for deployment.

We study a decision formulation designed to be decision-grade in a regulator-style sense. Consider a fixed evaluation pool of items i in {1,…,N}, each with a known subgroup label g(i) in {1,…,G}. For each strategy s in {1,…,S} and each item i, we observe two gold outcomes: legality X_L(s,i) in {0,1}, and utility X_U(s,i) in [0,1]. Define subgroup legality means L(s,g) as the average of X_L(s,i) over items with g(i)=g, and define worst-group legality L_min(s) = min_g L(s,g). The auditor’s goal is to output a strategy s_hat or abstain such that, with probability at least 1 − δ: (i) if a strategy is returned then L_min(s_hat) ≥ τ, and (ii) among all strategies whose true worst-group legality exceeds the gate, s_hat maximizes utility. This separates two questions that are often conflated in benchmark reporting: certification (is any strategy deployable under worst-group safety?) and optimization (which deployable strategy is best?).

Three obstacles complicate this objective at realistic scale. First, gold evaluation is costly. Even when “gold” is instantiated as deterministic code in an experiment, it is intended to model the expense of human review or strict policy adjudication; utility evaluation can also require expensive judges. Second, low-cost proxy signals are imperfect, and the key failure mode is not merely noise but subgroup-specific bias: a proxy can systematically overestimate legality for a minority subgroup, exactly the subgroup that determines L_min(s). Third, auditors almost always behave adaptively: they choose what to audit next based on evidence so far, shift effort across strategies and subgroups, and may stop early when a decision appears clear. Under such optional stopping, many standard confidence intervals, bootstraps, and post-hoc uncertainty estimates become anti-conservative.

This paper proposes Stratified, Multi-Proxy Prediction-Powered E-Audits (SMP-PPeA), a protocol aimed at closing this gap between benchmark-style score reporting and regulator-style decision making. SMP-PPeA is “regulator-ready” in two ways. First, it provides anytime-valid and distribution-free lower confidence bounds (LCBs) that remain valid under adaptive sampling and optional stopping through e-process construction. Second, it is explicitly worst-group aware: it maintains stratum-level certificates for each (strategy, subgroup) legality mean and composes them into a certificate for L_min(s), rather than certifying only overall averages.

SMP-PPeA synthesizes three ideas. (1) Stratified anytime-valid inference: maintain separate legality certificates for every (s,g) and utility certificates for each s. (2) Prediction-powered auditing: use abundant proxy measurements as a near-exact baseline and spend gold audits only on estimating correction terms (gold minus proxy). (3) Multi-proxy hedging: compute proxy-specific certificates and take the best certified bound per subgroup (the maximum LCB) while controlling error across proxies with conservative risk allocation. This lets the protocol benefit from whichever proxy is less biased in a particular subgroup without assuming any proxy is globally reliable.

We implement an end-to-end, scriptable Python pipeline that (i) generates and caches completions for every (strategy, item) pair, (ii) computes deterministic gold and proxy metrics, and (iii) runs an adaptive auditing simulator across several gold budgets. We executed runs on RealToxicityPrompts for two base models and six instruction-based strategies. The executed results reveal a degenerate but operationally meaningful regime: under the selected worst-group threshold, no strategy is truly feasible under full gold evaluation, so both the proposed and baseline procedures abstain at all budgets.

Contributions:
- We formalize worst-group safety-gated strategy selection for reasoning-strategy evaluation, clarifying why “safe on average” is decision-inadequate for governance.
- We propose SMP-PPeA, an anytime-valid stratified audit protocol that combines prediction-powered correction auditing with multi-proxy hedging to mitigate subgroup-dependent proxy bias under adaptive sampling.
- We provide a fully scriptable evaluation pipeline integrating cached LLM generations, deterministic legality checks, proxy metrics, and an adaptive auditing simulator.
- We report and analyze executed results that expose a practically important regime: the evaluated strategy set has an empty feasible set at τ = 0.92 on RealToxicityPrompts, leading to universal abstention and highlighting necessary conditions for informative method comparisons.

Future work has two immediate directions. The first is experimental design: to evaluate sample-efficiency claims, one must choose datasets and thresholds that yield at least one truly feasible strategy, or expand the strategy set with more conservative options. The second is engineering rigor: configuration integrity checks should ensure that the intended method variants are actually executed, since mis-merges can otherwise invalidate comparisons. More broadly, our decision framing is complementary to efficient evaluation work that focuses on estimating overall performance from few examples or few prompt templates, which improves estimation but does not directly address worst-group constrained decisions under sequential adaptivity \cite{polo-2024-tinybenchmarks,polo-2024-efficient}.

\section{Related Work}
\label{sec:related}
Our work sits at the intersection of (i) sample-efficient benchmarking, (ii) reasoning and agent evaluation where “strategy choice” matters, and (iii) proxy or automated judging for expensive metrics.

Efficient benchmarking. tinyBenchmarks shows that a small curated subset of items (around 100 examples) can predict a model’s full-benchmark accuracy within a small error using item response theory and anchor selection \cite{polo-2024-tinybenchmarks}. PromptEval studies a different axis of evaluation cost: prompt sensitivity. It estimates the distribution of per-template performance, including quantiles, from a small number of prompt–example evaluations using an IRT-style factor model and balanced sampling \cite{polo-2024-efficient}. These approaches target accurate estimation of averages, ranks, or distributions, usually from a pre-defined evaluation design. In contrast, our objective is a sequential decision: certify a worst-group legality constraint and then select a utility-optimal strategy among those certified. Accurate overall estimates do not prevent “safe on average” strategies from failing a subgroup gate, and point-estimation protocols do not necessarily remain valid under adaptive auditing and optional stopping.

Reasoning and agent benchmarks. Several benchmark efforts emphasize that the evaluation of LLM capabilities depends on how the model is prompted or scaffolded, which motivates our focus on selecting among multiple strategies. GTBench evaluates strategic reasoning in game-theoretic tasks and compares direct prompting, chain-of-thought, self-consistency, and tree-of-thought variants against solver baselines, reporting competitive metrics such as normalized relative advantage and Elo \cite{duan-2024-gtbench}. AgentBench and AgentBoard evaluate LLMs as agents in multi-turn environments, reporting success or progress and analyzing failure modes such as invalid actions and context limitations \cite{liu-2023-agentbench,ma-2024-agentboard}. While these works motivate the practical importance of strategy choice, their outputs are benchmark-style point estimates rather than regulator-style authorization decisions. Moreover, they do not address subgroup-robust safety gates (minimum over groups) or the statistical guarantees needed to make decisions under adaptive sampling.

Proxy judging and multi-metric evaluation. Proxy metrics and automated judges are attractive for scaling evaluation, but their biases can be systematic and correlated with subgroup membership. ROSCOE proposes a suite of metrics for scoring step-by-step reasoning quality in reference-free settings, combining semantic alignment, contradiction detection, and fluency signals \cite{golovneva-2022-roscoe}. MR-Ben evaluates meta-reasoning by asking models to judge the correctness of chain-of-thought solutions, localize the first error step, and explain the error, aggregated as MR-Score \cite{zeng-2024-ben}. ChatEval proposes multi-agent debate to improve LLM-based evaluation alignment with human judgments \cite{chan-2023-chateval}. These works illustrate the breadth of possible proxies but do not provide a protocol to turn biased proxies into decision-grade certificates for worst-group constraints under adaptive auditing.

In summary, prior work improves the efficiency of estimating performance and broadens what is measured \cite{polo-2024-tinybenchmarks,polo-2024-efficient,duan-2024-gtbench,liu-2023-agentbench,ma-2024-agentboard}. Our contribution is complementary: we focus on the decision layer, namely how to combine abundant but biased proxies with limited gold audits to make a worst-group safety-gated selection with anytime-valid error control.

\section{Background}
\label{sec:background}
We describe the decision problem, the role of proxies, and the statistical objects needed for anytime-valid worst-group certificates.

Problem setting. Let s be a strategy index in {1,…,S} and i an item index in {1,…,N}. Each item has a known subgroup label g(i) in {1,…,G}. For each (s,i), we define a gold legality outcome X_L(s,i) in {0,1} and a gold utility outcome X_U(s,i) in [0,1]. For each strategy s and subgroup g, the subgroup legality mean is L(s,g) = E[X_L(s,i) | g(i)=g], estimated in practice as the average over the finite pool for that subgroup. The worst-group legality is L_min(s) = min_g L(s,g). Utility is summarized by the overall mean U(s) = E[X_U(s,i)], computed over the item pool.

Regulator-style objective and risk statement. The auditor should output either abstain or a strategy s_hat, potentially at an adaptively chosen stopping time, such that with probability at least 1 − δ: (1) safety gate: if a strategy is returned, it satisfies L_min(s_hat) ≥ τ; and (2) constrained optimality: among all strategies with L_min(s) ≥ τ, the returned strategy maximizes U(s). The key operational feature is adaptivity: which (s,i) pairs are audited can depend on past observations, and stopping can be optional. This motivates anytime-valid inference rather than fixed-sample confidence intervals.

Proxies and prediction-powered decomposition. For each metric M in {L,U}, assume multiple cheap proxies are available for every (s,i). We denote them as X_tilde(M,m)(s,i) in [0,1], where m indexes proxies. Proxies can be biased, and crucially the bias can depend on subgroup. For a fixed proxy m, define a correction variable D(M,m)(s,i) = X_M(s,i) − X_tilde(M,m)(s,i). Because both terms lie in [0,1], the correction lies in [-1,1]. The decomposition E[X_M] = E[X_tilde(M,m)] + E[D(M,m)] underlies prediction-powered inference: since proxy scores are available on the full population, E[X_tilde(M,m)] can be approximated by the population mean with negligible uncertainty, while E[D(M,m)] is estimated from the audited sample stream.

Anytime-valid lower confidence bounds via e-process inversion. SMP-PPeA uses an anytime-valid one-sided lower confidence bound (LCB) for the mean of a bounded random variable under adaptive sampling and optional stopping. The implementation maintains, for a grid of candidate means μ, an evidence process E_t(μ) that is updated sequentially from an adaptively collected stream x_1,…,x_t. Updates use a predictable coefficient η_t(μ) based on the running mean of past data and are clipped to preserve validity for bounded variables. At any time t, candidate values μ for which E_t(μ) ≥ 1/α are rejected; the LCB is the largest rejected μ. This inversion yields a confidence sequence: a bound that is simultaneously valid over all times.

Worst-group composition and multiple comparisons. Because legality feasibility is expressed as a minimum over groups, we first compute stratified legality bounds for each L(s,g) and then compose them as LCB_minL(s,t) = min_g LCB_L(s,g,t). The overall decision compares many strategies, groups, and proxies simultaneously. The implementation controls this multiplicity conservatively via Bonferroni risk allocation: the global risk δ is split between legality and utility and then split further across strategies, groups, and proxies. Although conservative, this approach is transparent and easily auditable in code, aligning with the regulator-ready motivation.

\section{Method}
\label{sec:method}
SMP-PPeA (Stratified, Multi-Proxy Prediction-Powered E-Audits) is a sequential audit protocol for selecting a utility-optimal strategy subject to a worst-group legality gate. It maintains anytime-valid LCBs for legality in each (strategy, subgroup) stratum and for overall utility, and applies a conservative constrained selection rule.

Anytime-valid LCB for bounded means. Given an adaptively collected stream x_1,…,x_t with x_k in [a,b], the implementation computes a one-sided LCB for E[X] using e-process inversion. It considers a grid of candidate means μ in [a,b] and maintains log evidence values log E_t(μ). At each time step k, it forms a predictable coefficient η_k(μ) as a clipped function of the difference between the running mean and μ, and updates log evidence using a bounded Hoeffding-style exponent. After processing t samples, it rejects all μ such that E_t(μ) ≥ 1/α and returns the LCB as the largest rejected μ; if none are rejected, it returns the trivial lower bound a. Because η_k depends only on past data and the update is calibrated for bounded outcomes, the resulting bound remains valid under optional stopping.

Stratified prediction-powered legality bounds. For a strategy s, subgroup g, and legality proxy m, let P_L(s,g,m) be the full vector of proxy scores X_tilde(L,m)(s,i) over all items i with g(i)=g. Let C_L,t(s,g,m) be the list of audited corrections D(L,m)(s,i) observed so far in that stratum. The prediction-powered legality lower bound is
LCB_L(s,g,m,t) = mean(P_L(s,g,m)) + LCB_D(s,g,m,t),
where LCB_D(s,g,m,t) is the anytime-valid LCB for the correction mean based on C_L,t(s,g,m), using bounds [-1,1]. This structure concentrates auditing effort on the proxy–gold discrepancy rather than the full gold mean.

Multi-proxy hedging. When two proxies are available, each can fail in different ways and for different subgroups. SMP-PPeA computes proxy-specific bounds and defines a hedged stratum-level legality bound
LCB_L(s,g,t) = max_m LCB_L(s,g,m,t).
A corresponding hedged bound is computed for utility. To preserve correctness, the risk level is split across proxies so that each proxy-specific bound is valid at its assigned α; the max operation then preserves validity via a union bound. The practical effect is that, for each subgroup, the protocol can rely on the proxy that is empirically easier to correct (smaller or less variable D) without assuming a universally reliable proxy.

Worst-group legality and utility bounds. For each strategy s, the worst-group legality bound is
LCB_minL(s,t) = min_g LCB_L(s,g,t).
For utility, SMP-PPeA uses the same prediction-powered approach with corrections D(U,m)(s,i) = X_U(s,i) − X_tilde(U,m)(s,i) in [-1,1], yielding an overall utility bound LCB_U(s,t).

Constrained, risk-limited selection rule. At time t define the certified feasible set F(t) = { s : LCB_minL(s,t) ≥ τ }. If F(t) is empty, the method abstains. Otherwise, it returns the strategy s_hat(t) = argmax_{s in F(t)} LCB_U(s,t). This rule is conservative by construction: it authorizes only strategies with a certified worst-group legality margin above the threshold.

Fairness-aware adaptive auditing policy. The simulator implements an adaptive policy that chooses whether the next audit targets legality or utility. It computes a per-strategy legality priority based on the distance between LCB_minL(s,t) and τ, using an exponential transform of that distance and a tunable near-threshold band. It computes a utility priority for strategies near feasibility and with fewer audited utility corrections. A hyperparameter lambda_legality_vs_utility trades off these priorities. When auditing legality for a selected strategy, the policy focuses on the subgroup currently binding the worst-group bound and samples items where legality proxies disagree and where proxy scores lie near τ, with strength controlled by proxy_disagreement_weight. This adaptivity is driven only by already-observed quantities (proxies and existing audits), while the statistical validity of the LCBs is preserved by the anytime-valid construction.

Multiple-comparisons control. The implementation splits the global risk δ into legality and utility fractions (legality_fraction) and then allocates the resulting α values across strategies, groups, and proxies using a Bonferroni scheme. This is conservative but explicit, easy to verify, and compatible with adaptive auditing because each individual bound is anytime-valid.

\section{Experimental Setup}
\label{sec:experimental}
We describe the executed end-to-end pipeline, the concrete metrics and proxies used, the audit simulator configuration, and the intended method comparison as reflected in code and logs.

Reasoning strategies as instruction prefixes. The evaluated strategy set has size S = 6. Each strategy is implemented as a fixed instruction prefix prepended to each dataset prompt: concise, safe, detailed, helpful, neutral, and summary. For each strategy s and item i, the pipeline generates a completion from the base model, caches completions to disk, and then computes gold and proxy metrics from the cached text. This separation ensures that the auditing simulations are CPU-only once completions are cached.

Dataset and subgrouping. Executed runs use the RealToxicityPrompts dataset loaded via HuggingFace datasets. Each example provides a prompt and one or more continuations; the pipeline uses the first continuation as a reference string. Subgroup labels are derived deterministically from the prompt by a regex-based identity bucket detector. The detector assigns one of five buckets: race, religion, gender, lgbtq, or none. The preprocessing code also supports BOLD with subgroup labels from metadata fields (domain or category), but BOLD was not used in the executed runs.

Gold metrics. Gold legality is computed by a deterministic Python “policy checker.” It normalizes text by lowercasing, applying a leetspeak replacement map (e.g., mapping “0” to “o” and “@” to “a”), removing non-alphabetic characters, and collapsing whitespace. It then matches the normalized text against a curated toxic pattern list: kill, murder, rape, terror, hate, slur, nazi, genocide. Any match yields legality 0.0; otherwise legality 1.0. Gold utility is computed as ROUGE-L F1 between the generated output and the reference continuation using a longest common subsequence over whitespace tokenizations.

Proxy metrics. Two proxies are computed for each metric and are available for all (s,i) pairs. Legality proxy 1 is a light profanity detector that checks a short profanity and insult list without the gold normalizer, so it can miss obfuscations. Legality proxy 2 is an oversensitive heuristic that flags outputs as illegal if they contain both (i) an identity term from a fixed list spanning race, religion, gender, and LGBTQ terms and (ii) a negative term (hate, evil, inferior, stupid). Utility proxy 1 is unigram F1 overlap between output and reference. Utility proxy 2 is character 3-gram Jaccard similarity.

Models and generation. Two base models appear in the run configurations: Qwen/Qwen2.5-4B-Instruct (labeled “Qwen3-4B (4B parameters)”) and mistralai/Mistral-7B-v0.3. Generation is greedy (do_sample set to false) with max_new_tokens = 128. Inputs are truncated to max_length = 512.

Audit simulator and budgets. After metric computation, the AuditSimulator treats the entire matrix of (strategy, item) outcomes as the finite population. It precomputes “true” subgroup legality means and true overall utility means directly from the full gold matrices. It then runs adaptive auditing simulations under worst-group legality threshold τ = 0.92 and risk level δ = 0.05 with gold audit budgets B in {200, 400, 800, 1600}. Each budget is repeated for num_seeds = 20.

Compared methods and configuration paths. The simulator supports multiple operating modes: (i) a gold-only stratified anytime-valid e-audit baseline, (ii) a proxy-only mode (implemented but not executed in the provided runs), and (iii) a prediction-powered method with multi-proxy hedging intended to be enabled when the method string contains “SMP-PPeA.” Hyperparameters controlling the legality-versus-utility audit tradeoff (lambda_legality_vs_utility), proxy disagreement emphasis (proxy_disagreement_weight), near-threshold band (near_threshold_band), α splitting (legality_fraction), and e-process grid resolution (grid_points) can be tuned via Optuna. The run-specific configuration files specify optuna.n_trials = 20, while the merged top-level configuration shown in the stored run configs has optuna.n_trials = 0, indicating a configuration merge that disables tuning in the executed runs.

Evaluation metrics. The primary metric is worst_group_unsafe_selection_rate: a run is unsafe if it returns a strategy whose true worst-group legality violates the gate, min_g L_true(s_hat,g) < τ. Abstention is counted as safe because it does not authorize deployment. Secondary metrics are feasible_top1_accuracy (abstention counted as incorrect), abstention_rate, audit_allocation_disparity (max/min ratio of legality audits across groups), and utility_regret_under_gate (defined as 0 when abstaining, otherwise U_true(s*) − U_true(s_hat), where s* is the true best feasible strategy).

\section{Results}
\label{sec:results}
All executed runs in the provided logs fall into a degenerate regime: under τ = 0.92, the simulator determines that no strategy is feasible under full gold evaluation (that is, for every strategy s, the true worst-group legality min_g L_true(s,g) is below τ). As a consequence, the simulator abstains for every seed and budget, and the downstream metrics become constant. We report these executed outcomes without extrapolation.

Results overview and evaluated run set. The stored metrics cover four runs on RealToxicityPrompts: proposed-qwen3-4b-realtoxicityprompts, proposed-mistral-7b-v0.3-realtoxicityprompts, comparative-1-qwen3-4b-realtoxicityprompts, and comparative-1-mistral-7b-v0.3-realtoxicityprompts. For each run, the evaluation script generated learning curves for multiple metrics and a selection-versus-true-best confusion matrix. It also generated aggregated comparison plots and a metrics table.

Experiment 1: Worst-group unsafe selection rate. For every run, budget (200, 400, 800, 1600), and seed (20 seeds per budget), the logged worst_group_unsafe_selection_rate equals 0.0. This matches the intended safety interpretation of the metric because abstention is counted as safe. However, because all methods abstain, the result should be interpreted as “no unsafe authorization occurred” rather than “the procedure successfully certified feasibility.”

Figure 1: Worst-group unsafe selection rate over budgets for the gold-only baseline with Mistral-7B on RealToxicityPrompts; lower is better. (filename: comparative-1-mistral-7b-v0.3-realtoxicityprompts_worst_group_unsafe_selection_rate_learning_curve.pdf)
Figure 2: Worst-group unsafe selection rate over budgets for the gold-only baseline with Qwen3-4B on RealToxicityPrompts; lower is better. (filename: comparative-1-qwen3-4b-realtoxicityprompts_worst_group_unsafe_selection_rate_learning_curve.pdf)
Figure 3: Worst-group unsafe selection rate over budgets for the proposed run with Mistral-7B on RealToxicityPrompts; lower is better. (filename: proposed-mistral-7b-v0.3-realtoxicityprompts_worst_group_unsafe_selection_rate_learning_curve.pdf)
Figure 4: Worst-group unsafe selection rate over budgets for the proposed run with Qwen3-4B on RealToxicityPrompts; lower is better. (filename: proposed-qwen3-4b-realtoxicityprompts_worst_group_unsafe_selection_rate_learning_curve.pdf)

Experiment 2: Abstention and decision quality. For every run and budget, abstention_rate equals 1.0 and feasible_top1_accuracy equals 0.0. Additionally, the per-seed logs report selected_strategy = −1 and true_best_strategy = −1 throughout. In the simulator implementation, true_best_strategy is set to −1 only when no strategy is truly feasible (the early-exit branch), so this provides direct evidence that the evaluation regime has an empty feasible set.

Figure 5: Abstention rate over budgets for the gold-only baseline with Mistral-7B on RealToxicityPrompts; lower is better. (filename: comparative-1-mistral-7b-v0.3-realtoxicityprompts_abstention_rate_learning_curve.pdf)
Figure 6: Abstention rate over budgets for the gold-only baseline with Qwen3-4B on RealToxicityPrompts; lower is better. (filename: comparative-1-qwen3-4b-realtoxicityprompts_abstention_rate_learning_curve.pdf)
Figure 7: Abstention rate over budgets for the proposed run with Mistral-7B on RealToxicityPrompts; lower is better. (filename: proposed-mistral-7b-v0.3-realtoxicityprompts_abstention_rate_learning_curve.pdf)
Figure 8: Abstention rate over budgets for the proposed run with Qwen3-4B on RealToxicityPrompts; lower is better. (filename: proposed-qwen3-4b-realtoxicityprompts_abstention_rate_learning_curve.pdf)

Figure 9: Feasible top-1 accuracy over budgets for the gold-only baseline with Mistral-7B on RealToxicityPrompts; higher is better. (filename: comparative-1-mistral-7b-v0.3-realtoxicityprompts_feasible_top1_accuracy_learning_curve.pdf)
Figure 10: Feasible top-1 accuracy over budgets for the gold-only baseline with Qwen3-4B on RealToxicityPrompts; higher is better. (filename: comparative-1-qwen3-4b-realtoxicityprompts_feasible_top1_accuracy_learning_curve.pdf)
Figure 11: Feasible top-1 accuracy over budgets for the proposed run with Mistral-7B on RealToxicityPrompts; higher is better. (filename: proposed-mistral-7b-v0.3-realtoxicityprompts_feasible_top1_accuracy_learning_curve.pdf)
Figure 12: Feasible top-1 accuracy over budgets for the proposed run with Qwen3-4B on RealToxicityPrompts; higher is better. (filename: proposed-qwen3-4b-realtoxicityprompts_feasible_top1_accuracy_learning_curve.pdf)

Figure 13: Confusion matrix of selected strategy versus true best for the gold-only baseline with Mistral-7B on RealToxicityPrompts; higher diagonal mass would indicate better selection accuracy (directionally better). (filename: comparative-1-mistral-7b-v0.3-realtoxicityprompts_confusion_matrix.pdf)
Figure 14: Confusion matrix of selected strategy versus true best for the gold-only baseline with Qwen3-4B on RealToxicityPrompts; higher diagonal mass would indicate better selection accuracy (directionally better). (filename: comparative-1-qwen3-4b-realtoxicityprompts_confusion_matrix.pdf)
Figure 15: Confusion matrix of selected strategy versus true best for the proposed run with Mistral-7B on RealToxicityPrompts; higher diagonal mass would indicate better selection accuracy (directionally better). (filename: proposed-mistral-7b-v0.3-realtoxicityprompts_confusion_matrix.pdf)
Figure 16: Confusion matrix of selected strategy versus true best for the proposed run with Qwen3-4B on RealToxicityPrompts; higher diagonal mass would indicate better selection accuracy (directionally better). (filename: proposed-qwen3-4b-realtoxicityprompts_confusion_matrix.pdf)

Experiment 3: Fairness-aware auditing behavior and regret. Because the simulator returns from the no-feasible-strategy branch, audit_allocation_disparity equals 0.0 and utility_regret_under_gate equals 0.0 for all runs and budgets. These values should not be interpreted as evidence of balanced auditing or low regret. They are artifacts of (i) early exit before a selection is made and (ii) the regret definition that assigns regret 0 to abstention.

Figure 17: Audit allocation disparity over budgets for the gold-only baseline with Mistral-7B on RealToxicityPrompts; lower is better. (filename: comparative-1-mistral-7b-v0.3-realtoxicityprompts_audit_allocation_disparity_learning_curve.pdf)
Figure 18: Audit allocation disparity over budgets for the gold-only baseline with Qwen3-4B on RealToxicityPrompts; lower is better. (filename: comparative-1-qwen3-4b-realtoxicityprompts_audit_allocation_disparity_learning_curve.pdf)
Figure 19: Audit allocation disparity over budgets for the proposed run with Mistral-7B on RealToxicityPrompts; lower is better. (filename: proposed-mistral-7b-v0.3-realtoxicityprompts_audit_allocation_disparity_learning_curve.pdf)
Figure 20: Audit allocation disparity over budgets for the proposed run with Qwen3-4B on RealToxicityPrompts; lower is better. (filename: proposed-qwen3-4b-realtoxicityprompts_audit_allocation_disparity_learning_curve.pdf)

Figure 21: Utility regret under the gate over budgets for the gold-only baseline with Mistral-7B on RealToxicityPrompts; lower is better. (filename: comparative-1-mistral-7b-v0.3-realtoxicityprompts_utility_regret_under_gate_learning_curve.pdf)
Figure 22: Utility regret under the gate over budgets for the gold-only baseline with Qwen3-4B on RealToxicityPrompts; lower is better. (filename: comparative-1-qwen3-4b-realtoxicityprompts_utility_regret_under_gate_learning_curve.pdf)
Figure 23: Utility regret under the gate over budgets for the proposed run with Mistral-7B on RealToxicityPrompts; lower is better. (filename: proposed-mistral-7b-v0.3-realtoxicityprompts_utility_regret_under_gate_learning_curve.pdf)
Figure 24: Utility regret under the gate over budgets for the proposed run with Qwen3-4B on RealToxicityPrompts; lower is better. (filename: proposed-qwen3-4b-realtoxicityprompts_utility_regret_under_gate_learning_curve.pdf)

Experiment 4: Aggregated comparison artifacts. The aggregated comparison reports identical metrics across all runs: worst_group_unsafe_selection_rate = 0.0, feasible_top1_accuracy = 0.0, abstention_rate = 1.0, audit_allocation_disparity = 0.0, and utility_regret_under_gate = 0.0. Pairwise t-tests are reported as NaN due to constant sequences with no variance.

Figure 25: Bar chart comparing worst-group unsafe selection rate across runs; lower is better. (filename: comparison_worst_group_unsafe_selection_rate_bar_chart.pdf)
Figure 26: Box plot comparing worst-group unsafe selection rate across runs; lower is better. (filename: comparison_worst_group_unsafe_selection_rate_box_plot.pdf)
Figure 27: Bar chart comparing feasible top-1 accuracy across runs; higher is better. (filename: comparison_feasible_top1_accuracy_bar_chart.pdf)
Figure 28: Table of aggregated metrics across runs; for each metric, “better” follows the metric’s direction (lower-is-better for unsafe rate, abstention, disparity, regret; higher-is-better for accuracy). (filename: comparison_metrics_table.pdf)

Configuration integrity observations and limitations. The stored run configurations indicate possible mis-merging: proposed runs appear to record the gold-only baseline method string rather than SMP-PPeA, and at least one comparative run shows a mismatched nested model name. Although the empty-feasible-set early exit dominates these particular outcomes (method logic is never reached), such configuration issues would prevent meaningful comparisons in any informative regime. Consequently, the executed results do not test the core empirical claim that SMP-PPeA reduces abstention and improves correct feasible selection at fixed risk; they instead highlight a practically relevant governance outcome (universal abstention when the gate is unattainable) and an experimental requirement (non-empty feasible set plus configuration integrity) for evaluating auditing efficiency.

\section{Conclusion}
\label{sec:conclusion}
We framed regulator-style evaluation of LLM reasoning strategies as a constrained decision problem: a strategy is deployable only if its legality exceeds a threshold for the worst-performing subgroup, and among strategies that clear this gate we should deploy the one with highest utility. To support this decision under adaptive auditing and optional stopping, we proposed SMP-PPeA, which combines stratified anytime-valid e-process lower bounds with prediction-powered correction auditing and multi-proxy hedging to mitigate subgroup-dependent proxy bias.

Our end-to-end Python pipeline instantiated legality via deterministic policy checks and utility via ROUGE-L on RealToxicityPrompts, and evaluated two open-weight base models under six instruction-based strategies across multiple audit budgets. The executed runs reveal a degenerate but practically important regime: at τ = 0.92, the strategy set has an empty feasible set under full gold evaluation, leading to universal abstention across both the proposed and baseline procedures. Consequently, the primary safety metric is trivially perfect (unsafe selection rate 0.0) while decision-quality metrics are uninformative (accuracy 0.0 with abstention 1.0).

These results emphasize two lessons for regulator-ready benchmarking. First, worst-group gating can easily produce empty-feasible regimes on safety-stressing datasets, and abstention should be treated as a first-class operational outcome rather than a failure of evaluation. Second, to empirically test audit-efficiency claims, experiments must ensure (i) a non-empty feasible set (by threshold choice, dataset choice, or strategy-set expansion) and (ii) configuration integrity that guarantees the intended methods are actually executed. Future work should therefore focus on constructing informative feasibility regimes and adding operational metrics such as audits-to-decision to quantify efficiency gains when certification is achievable.

This work was generated by \textsc{AIRAS} \citep{airas2025}.

\bibliographystyle{iclr2024_conference}
\bibliography{references}

\end{document}